% basic nastaveni
\documentclass[12pt]{article}
\setlength{\parindent}{0pt}
\setlength{\parskip}{1em}
\usepackage[utf8]{inputenc}
\usepackage[czech]{babel}
\usepackage[T1]{fontenc}
\usepackage{graphics}
\usepackage{caption}
\usepackage{hyperref}
\usepackage{xcolor}
\usepackage{float}

% matematicke package
\usepackage{amsmath}
\usepackage{amsfonts}
\usepackage{amssymb}
\usepackage{geometry}
\geometry{margin=2.5cm}
\usepackage{pgfplots}
\pgfplotsset{compat=1.18}
\usepackage{mathrsfs}

\title{Maturitní otázka číslo 4: Algebraické rovnice, nerovnice a jejich soustavy}
\author{}
\date{}

\begin{document}
\maketitle

\subsection*{Definice ronvice, nerovnice, definiční obor a množina kořenů}
Mějme dva výrazy $L(x), P(x)$ s proměnnou $x$. Určují se hodnoty z číselného oboru $M$ pro něž platí
rovnost u rovnic nebo nerovnost u nerovnic. $L(x)$ se nazývá levá strana rovnice(nerovnice) a $R(x)$
pravá. V některých případech může být neznámá pouze v jednom z výrazů. Proměnná $x$ se nazývá neznámá.Hodnoty neznámé pro něž je rovnice(nerovnice) splněna se nazývá kořenem. Číselný obor $M$ je základní
množina čísel, ve které hledáme kořeny rovnice. Podmnožina číselného oboru, v níž jsou definované
oba výrazy $L(x), P(x)$ se nazývá definicční obor a značí se $D$. Množina všech kořenů rovnic a
nerovnic se nazývá množina všech kořenů a značí se $K$.
\subsection*{Úpravy rovnic a nerovnic}
\subsubsection*{Úprávy rovnic}
Úpravy rovnic se dělí na ekvivalentní a neekvivalentní. Ekvivalentní upráva rovnice nezmění rovnost
rovnice. Tedy po úpravě výrazu $L(x)$ na $L_1(x)$ a $R(x)$ na $R_1(x)$, pořád platí $L_1(x) = P_1(x)$.
Neekvivalentní úpravy vytvoří také novou rovnost, ale může se změnit množina kořenů, jako například
přidáním falešného kořenu nebo odebráním realného kořenu. Tyto úpravy se mohou provádět pokud si
uvědomujeme jejich rizika a předem uděláme kroky k zaručení ekvivalence.

První operací je prohození stran rovnice. Jelikož se obě strany rovnají, mohu je bezpečně prohodit.

Další operací je u rovnic sčítání a odčítání. Přičteš nebo odečteš stejné číslo nebo výraz
k oboum stranám rovnice.
\[
\begin{aligned}
    x - 5 &= -2 \\
    x &= 3
\end{aligned}
\]
Další operací je násobení a dělení. Při násobení a dělení se násobí a dělí všechny členy
výrazu obou stran nějakým číslem nebo výrazem. Dá se to představit jako bychom dali celý výraz do
závorek a provedli tuto operaci na této závorce.

\[
\begin{aligned}
    \frac{3x}{7} &= 9 \\
    \frac{x}{7} &= 3 \\
    x &= 21
\end{aligned}
\]
Násobení nulou je neekvivalentní úprava. Dělení nulou není definované. Pozor si musíme dát při
násobení a dělení neznámými.
\[
\begin{aligned}
    x^2 &= 5x \\
    x &= 5
\end{aligned}
\]
Po vydělení celé rovnice $x$ máme jedno řešení, ovšem nula je také řešením.

Další operací je umocňování a odmocňování, které se také provádí pro celé výrazy obou stran rovnice.
Tyto operace ovšem nejsou vždy ekvivalentní. Umocňování sudým číslem může převrátit znaménka. Pokud
chceme aby umocňování sudým číslem bylo ekvivalentní, musí být obě strany rovnice nezáporné.
\[
\begin{aligned}
    x &= 3 \\
    x^2 &= 9
\end{aligned}
\]
Toto umocnění vytvořilo falešný kořen $-3$.

Odmocnění je vždy ekvivalentní pokud odmocňujeme lichým číslem. Pokud odmocňujeme sudým číslem,
musíme dát část výrazu s neznámou do absolutních hodnot.
\[
\begin{aligned}
    4x^2 &= 64 \\
    \sqrt{4x^2} &= 8 \\
    2|x| &= 8 \\
    |x| &= 4
\end{aligned}
\]
Pokud bychom nedali $x$ do absolutních hodnot, ztratily bychom řešení $-4$.
\subsubsection*{Úpravy nerovnic}
Pokud chcme u nerovnic prohodit strany, musíme také prohodit znaménko nerovnosti, jinak se nejedná
o ekvivalentní úpravu.

Přičítání a odečítání funguje stejně jako u rovnic a jedná se o ekvivalentní úpravu.

Pokud násobíme nebo dělíme nerovnici kladným číslem, jde o ekvivalentní úpravu. Pokud to je ovšem
záporným číslem, musíme prohodit znaménko nerovnosti. Násobení a dělení nulou funguje stejně jako
u rovnic. Pozor si musíme dát na násobení a dělení neznámou či výrazem s neznámou. V ten moment
se nerovnice rozdělí na dvě a to na první ve které je neznámá nebo výraz kladný a na druhou kde
je záporný. Sjednocení kořenů těchto dvou nerovnic je pak finální řešení.

Umocňování a odmocňování se většinou u nerovnic neprovádí, jelikož je ošetřování nerovnice poměrně
zdlouhavé.
\subsubsection*{Zkoužka}
Pokud nalezneme nějaké řešení rovnice nebo nerovnice, je dobré zjistit, zda je opravdu správné.
Mějme rovnici $L(x) = P(x)$ a řešení $x_0$. Po dosazení $x_0$ do levé i pravé strany zjišťujeme, zda
rovnost pořád platí.
\subsection*{Polynomy, nulové body a grafické řešení rovnic}
Polynom je výraz ve tvaru $a_nx^n + a_{n-1}x^{n-1} + \dots + a_1x + a_0$. Čísla $a_n, a_{n-1}, \dots,
a_0$ se nazývají koeficienty.

Pokud máme nějakou rovnici, lze jí vždy upravit tak, že budeme mít na pravé straně nulu. Pokud
bychom si představili, že výraz na levé straně je funkce, řešením této rovnice budou body protínající
osu $x$ neboli nulové body. To je jedna z možností jak řešit rovnice graficky. Další možností je
si představit, že obě strany rovnice jsou vlastně funkce a Řešením jsou průsečíky těchto dvou funkcí.

U nerovností je grafické řešení podobné. Také si přestavíme že výrazy na obou stranách jsou vlastně
funkce a pozorujeme, kdy je jedna funkce menší než druhá. Narozdíl od rovnic bude výsledkem interval.
\subsection*{Řešení lineárních a kvadratických rovnic a nerovnic}
Řešení lineárních rovnic a nerovnic je pouze využívaní možných úprav k izolování neznámé.
Může občas vyjít speciální případy a to nekonečně mnoho řešení a žádné řešení. Pokud po úpravě vyjde
$0 = 0$ rovnice má nekonečno mnoho řešení. Pokud vyjde $0 = k, k \in mathbb{R}$, rovnice nemá žádné
řešení. Příkladem můžou být rovnice $x = x$ pro nekonečně mnoho řešení a rovnice $x = x+1$ pro žádné
řešení.

Pro řešení kvadratických rovnic a nerovnic se nejdříve převede vše doleva aby na pravé straně zbyla
nula. Následně se využije jedna z metod a to vytýkání, Vietovy vztahy, doplnění na čtverec nebo pomocí
diskriminantu. Kvadratické rovnice mají v oboru realných čísel dvě, jedno nebo žádné řešení. Mějme
levou stranu $ax^2 + bx + c$.

Vytíkání se využívá, pokud chybí absolutní člen, tedy člen $c$. Vytknutím $ax$ dostaneme
$ax (x + \frac {b}{a})$ což nám už dává nulové body a to $x_1 = 0, x_2 = - \frac{b}{a}$.

Vietovy vztahy nám říkají, že $x_1 + x_2 = -\frac{b}{a}, x_1 \cdot x_2 = \frac{c}{a}$. Vzorce se
nejlépe využívá pokud je kvadratický člen roven jedné.

Doplnění na čtverec hledá vzorec $(a \pm b^2)$. V polynomu se nachází začátek vzorce $x^2 \pm 2xy +
y^2$ a mi si k prvním dovu částem uměle dotvoříme tento vzorec. Z doplnění na čtverec vznikl známý
vzorec pro diskriminant tím, že se použije doplnění na čvterec na základní předpis funkce.

Zjištění kořenů pomocí diskriminantu se využívá pokud jsou kořeny rovnice na které nejde lehce přijít.
Diskriminat získáme vzorcem $D = b^2 - 4ac$. Pokud je $D$ kladné, rovnice má dvě řešení, pokud je
roven nule, rovnice má jedno řešení a je to dvojnásobný reálný kořen. Pokud je záporný, rovnice nemá
v oboru reálných čísel řešení. Kořeny získáme ze vzorce $x_{1,2} = \frac{-b \pm \sqrt{D}}{2a}$.

Ukažme si všechny možnosti řešení kromě vytknutí v praxi na rovnici $x^2 - 6x + 8 = 0$
\[
\begin{aligned}
    (x - 4)(x - 2) &= 0 \\
    4 + 2 = 6, -4 \cdot -2 &= 8 \\
    \hline \\
    x^2 -6x +9 -9 +8 &= 0 \\
    (x - 3)^2 -9 + 8 &= 0 \\
    (x - 3)^2 &= 1 \\
    |x - 3| &= 1 \\
    \hline \\
    \frac{6 \pm \sqrt{36 - 4 \cdot 1 \cdot 8}}{2} &= x_{1,2} \\
    3 \pm 1 &= x_{1,2}
\end{aligned}
\]
Řešení se může zapsat dvěmi způsoby a to $K = \{2, 4\}$ nebo $x \in \{2, 4\}$

Pokud bychom měli nerovnici $x^2 - 6x + 8 < 0$ nejdříve bychom zjistili nulové body uplně stejným
způsobem, čímž bychom si rozdělili definiční obor na tři části $(-\infty, 2), (2, 4), (4, \infty)$.
Následně bychom z každého intervaly vybrali číslo a dosadili ho do nerovnice.
\[
\begin{aligned}
    8 &< 0 \\
    \hline \\
    3^2 - 18 + 8 &< 0 \\
    -1 &< 0 \\
    \hline \\
    10^2 - 60 + 8 &< 0 \\
    48 &< 0
\end{aligned}
\]
Z toho nám vychází že řešení je $K = (2, 4)$. Pokud by místo ostré nerovnosti byla v zadání neostrá,
byli by nulové body součástí řešení, tedy $K = \langle 2, 4 \rangle$.

\subsection*{Rovnice a nerovnice s parametrem}
Rovnice s parametrem se řeší tím, že si necháme jenom proměnou na jedné stranně a zbytek dáme
na druhou, následně vytvoříme tabulku s tím jaké hodnoty získáme pro jaký parametr. Mějme rovnici
\[
\begin{aligned}
    a^2x - a &= x + 1 \\
    x(a^2 - 1) &= a + 1 \\
    x &= \frac{a+1}{(a-1)(a+1)} \\
    x &= \frac{1}{a-1}
\end{aligned}
\]
A následně se všechna řešení dají do tabulky. Pro $a \ne 1, -1, x = \frac{1}{a-1}$. Pokud dosadíme
$a = 1$ do původní rovnice, dostaneme $0 = 2$, tedy $x$ nemá řešení. Pro $a = -1$ získáme $0 = 0$,
tedy $x$ je jakékoliv realné číslo.

\begin{center}
\begin{tabular}{|l|l|}
\hline
$a$ & $x$ \\
\hline
$a \ne 1, -1$ & $x = \frac{1}{a-1}$ \\
\hline
$a = 1$ & $x \in \emptyset$ \\
\hline
$a = -1$ & $x \in \mathbb{R}$ \\
\hline
\end{tabular}
\end{center}

Pokud bychom měli kvadratickou rovnici $x^2 -ax + a + 3 = 0$, rozdělili bychom řešení podle toho jak
vychází diskriminant.
\[
\begin{aligned}
    D &= (-a)^2 - 4 \cdot 1 \cdot (a+3) \\
    D &= (a-6)(a+2)
\end{aligned}
\]
Diskriminant je záporný v rozmezí od $(-2, 6)$, takže rovnice nemá řešení. Diskriminant je roven nule
pro $a = 6, a = -2$. Tyto hodnoty můžeme dosadit do zbytku vzorce pro řešení $\frac{-b}{2a}$ a pro
$a = 6$ dostaneme $3$ a pro $a = -2$ dostaneme $-1$. Diskriminant je klady v useku $(-\infty, -2) 
\cup (6, \infty)$ a rovnice má řešení $x_{1,2} = \frac{a \pm \sqrt{a^2 - 4a -12}}{2}$.
\begin{center}
\renewcommand{\arraystretch}{1.5} 
\begin{tabular}{|l|l|}
    \hline
    $m$ & $x$ \\
    \hline
    $m \in (-2; 6)$ & $x \in \emptyset$ \\
    \hline
    $m = 6$ & $x = \{3\}$ \\
    \hline
    $m = -2$ & $x = \{-1\}$ \\
    \hline
    $m \in (-\infty; -2) \cup (6; \infty)$ & $x = \left\{ \frac{m \pm \sqrt{m^2 - 4m - 12}}{2} \right\}$ \\
    \hline
\end{tabular}
\end{center}
S nerovniemi to funguje uplně stejně, ale pracuje se většinou s velkým množstvím štěpení nerovnice.
\end{document}
